\section{Area and modular operators in the sectors \texorpdfstring{\(90\)}{90} and \texorpdfstring{\(120\)}{120}}

In each of the two families of
\cref{rm:eq:modular-boundary-counts} there are eighteen qutrit channels. For
\(m\in\{90,120\}\), define
\begin{equation}
\label{rm:eq:modular-number-operators}
N_m=\sum_{e\in\partial A_m}N_e,
\qquad \Spec(N_e)=\{0,1,2\}.
\end{equation}
The Barbero--Immirzi output is used here as a prior result, without
rederiving it. With the area normalization constructed in
\cref{rm:eq:modular-theta-clock-a0}, one obtains
\begin{equation}
\label{rm:eq:modular-area-operator}
\boxed{
\frac{\widehat{\mathcal A}_{\rm phys}}{\ell_P^2}
=\gammaH a_0(N_{90}+N_{120}),}
\qquad
\gammaH a_0=0.0016755940749224991.
\end{equation}

\begin{theorem}[Finite spectrum of the area quantum]
\label{rm:thm:modular-area-spectrum}
On the thirty-six links of the cut,
\begin{equation}
\label{rm:eq:modular-total-number-spectrum}
\Spec(N_{90}+N_{120})=\{0,1,\ldots,72\}.
\end{equation}
The multiplicity of the eigenvalue \(n\) is the coefficient of \(z^n\) in
\((1+z+z^2)^{36}\). Consequently,
\begin{equation}
\label{rm:eq:modular-area-spectrum}
\boxed{
\Spec\!\left(\frac{\widehat{\mathcal A}_{\rm phys}}{\ell_P^2}\right)
=\{n\gammaH a_0:0\leq n\leq72\}.}
\end{equation}
The compatible cuts prolong the family
\(\mathcal A_n^{\rm prot}=n\gammaH a_0\ell_P^2\).
\end{theorem}

\begin{proof}
Each operator \(N_e\) has a spectral generating function
\(1+z+z^2\). The thirty-six factors commute and act on distinct
tensor factors; therefore, the generating function of the total is
\((1+z+z^2)^{36}\), whose degrees range over all integers from zero to
seventy-two. Multiplication by the positive scalar
\(\gammaH a_0\ell_P^2\) gives \cref{rm:eq:modular-area-spectrum}.
\end{proof}

Two quantization procedures must be distinguished. The cut law \cref{rm:thm:modular-triadic-cut} counts eighty-one unit capacities and determines \(81\log3\); the operator \cref{rm:eq:modular-area-operator} quantizes occupations of thirty-six response links, up to seventy-two number trits. They share the same boundary sheet, but not the same observable.

The torsion modular produces
\begin{equation}
\label{rm:eq:modular-delta-value}
\Dmod
=0.0001111970500599773249414566131933856\ldots,
\end{equation}
and the boundary Hamiltonian
\begin{equation}
\label{rm:eq:modular-hamiltonian}
\boxed{
\HHM^{(A)}=-\log\rho_A
=cI+\Dmod(90N_{90}+120N_{120}).}
\end{equation}
The qualifier \emph{holographic} does not designate an analogy: the
incidence \(6\to8\) is transported to the boundary sets
\(\partial A_{90}\) and \(\partial A_{120}\), and the reduced state resides
in the algebra generated by its number operators.

\begin{theorem}[Joint reconstruction of geometry and memory]
\label{rm:thm:modular-area-memory}
Let
\begin{equation}
\label{rm:eq:modular-ND}
N=N_{90}+N_{120},
\qquad
D=90N_{90}+120N_{120}.
\end{equation}
Then
\begin{equation}
\label{rm:eq:modular-inverse-channels}
\boxed{
N_{90}=4N-\frac{D}{30},
\qquad
N_{120}=\frac{D}{30}-3N.}
\end{equation}
In particular,
\([\widehat{\mathcal A}_{\rm phys},\HHM^{(A)}]=0\), and the pair
\((\widehat{\mathcal A}_{\rm phys},\HHM^{(A)}-cI)\) reconstructs the two
occupation sectors.
\end{theorem}

\begin{proof}
The operators \(N_{90}\) and \(N_{120}\) act on distinct factors
and commute. The linear map carrying \((N_{90},N_{120})\) to \((N,D)\)
has matrix
\[
\begin{pmatrix}1&1\\90&120\end{pmatrix},
\qquad \det=30.
\]
Its inverse is exactly \cref{rm:eq:modular-inverse-channels}. Since
\cref{rm:eq:modular-area-operator,rm:eq:modular-hamiltonian} are linear
combinations of the same number operators, the commutator vanishes.
\end{proof}

The operator interpretation is precise: the area observable preserves the total number of excitations, while the modular Hamiltonian preserves their sector of origin. Neither observable alone retains both coordinates; the joint pair does determine them.

\section{Confluence of area, information and energy}

The angular scale \(\theta_{\rm clk}\) in
\cref{rm:eq:modular-theta-clock-a0} and the chronological temperature
\(\Theta_{\rm clk}\) are distinct objects; the use of an uppercase letter is part
of the typing. The thermodynamic realization in \cref{rm:ch:metrology} proves
\begin{equation}
\label{rm:eq:modular-planck-trit-bridge}
E_P=k_B\Theta_{\rm clk}\log3,
\end{equation}
whereas the areal realization constructs the elementary increment
\begin{equation}
\label{rm:eq:modular-area-quantum}
\delta\mathcal A
=\gammaH a_0\ell_P^2.
\end{equation}
They are not equated: \cref{rm:eq:modular-planck-trit-bridge} fixes the informational quantum of energy, and \cref{rm:eq:modular-area-quantum} the geometric quantum of area. Their quotient defines the typed tension of conversion
\begin{equation}
\label{rm:eq:modular-information-area-tension}
\boxed{
\mathcal T_{I/A}
=\frac{E_P}{\delta\mathcal A}
=\frac{k_B\Theta_{\rm clk}\log3}
{\gammaH a_0\ell_P^2}.}
\end{equation}
The expression introduces no additional constant: it determines
the normalizations that must be preserved in the transformation
of information into geometry.


In the minimal cut of eighty-one links, the physical information capacity and the energy of one trit per link are
\begin{equation}
\label{rm:eq:modular-cut-info-energy}
S_{\rm cut}^{\rm phys}=81k_B\log3,
\qquad
E_{\rm cut}^{(1)}=81E_P
=\Theta_{\rm clk}S_{\rm cut}^{\rm phys}.
\end{equation}
The final equality is an equation of state of this realization, not a universal
equipartition law. In particular, the Gibbs state of the thirty-six
boundary links has a different entropy and is treated next.

This separation determines the structural content of the confluence:
Barbero--Immirzi
orients the area spectrum; \(\log3\) measures the irreducible decision of the
TRIT; \(\Theta_{\rm clk}\) converts that information into energy; and
\(\ell_P^2\) connects the area quantum to the gravitational family.
Gravity is described as a confluence of functional realizations of a
common antecedent, and not as an imposed equality among their values.
